\documentclass[aspectratio=169,10pt]{beamer}

% Shared Beamer setup for Fall 2026 course slides.
\mode<presentation>{
  \usetheme{Madrid}
  \usecolortheme{default}
}

\definecolor{CalPolyGreen}{HTML}{154734}
\definecolor{CalPolyGold}{HTML}{BD8B13}
\definecolor{DeckInk}{HTML}{1F2933}
\definecolor{DeckMuted}{HTML}{5F6B7A}

\setbeamercolor{structure}{fg=CalPolyGreen}
\setbeamercolor{palette primary}{bg=CalPolyGreen,fg=white}
\setbeamercolor{palette secondary}{bg=DeckInk,fg=white}
\setbeamercolor{palette tertiary}{bg=CalPolyGold,fg=DeckInk}
\setbeamercolor{frametitle}{bg=CalPolyGreen,fg=white}
\setbeamercolor{title}{fg=CalPolyGreen}
\setbeamercolor{normal text}{fg=DeckInk,bg=white}
\setbeamercolor{block title}{bg=CalPolyGreen,fg=white}
\setbeamercolor{block body}{bg=black!3,fg=DeckInk}

\setbeamertemplate{navigation symbols}{}
\setbeamertemplate{footline}[frame number]
\setbeamertemplate{itemize item}{\color{CalPolyGold}$\blacktriangleright$}
\setbeamertemplate{itemize subitem}{\color{CalPolyGreen}$\bullet$}

\newcommand{\CourseNumber}{Course}
\newcommand{\CourseTitle}{Course Title}
\newcommand{\LectureNumber}{0}
\newcommand{\LectureTitle}{Lecture Title}
\newcommand{\LectureDate}{Date}
\newcommand{\CourseUnit}{Unit}

\newcommand{\coursenumber}[1]{\renewcommand{\CourseNumber}{#1}}
\newcommand{\coursetitle}[1]{\renewcommand{\CourseTitle}{#1}}
\newcommand{\lecturenumber}[1]{\renewcommand{\LectureNumber}{#1}}
\newcommand{\lecturetitle}[1]{\renewcommand{\LectureTitle}{#1}}
\newcommand{\lecturedate}[1]{\renewcommand{\LectureDate}{#1}}
\newcommand{\courseunit}[1]{\renewcommand{\CourseUnit}{#1}}

\author{Kirk Duran}
\institute{California Polytechnic State University, San Luis Obispo}
\date{\LectureDate}

\newcommand{\makedecktitle}{
  \begin{frame}[plain]
    \vspace{0.25in}
    {\usebeamerfont{title}\color{CalPolyGreen}\LectureTitle\par}
    \vspace{0.18in}
    {\large \CourseNumber{}: \CourseTitle\par}
    \vspace{0.08in}
    {\large Lecture \LectureNumber{} -- \LectureDate\par}
    \vfill
    {\small Kirk Duran \textbar{} Cal Poly, San Luis Obispo\par}
  \end{frame}
}

\newcommand{\placeholdernote}{
  \begin{block}{Placeholder}
    Replace these bullets with the final lecture material, examples, and in-class activities.
  \end{block}
}

\AtBeginSection[]{
  \begin{frame}{Roadmap}
    \tableofcontents[currentsection]
  \end{frame}
}


\usepackage{tikz}

\graphicspath{{../assets/lecture-10/}}

% If an image file is not yet available, skip the visual slot.
\newcommand{\lectureimage}[2][1.75in]{%
  \IfFileExists{../assets/lecture-10/#2}{%
    \includegraphics[width=\linewidth,height=#1,keepaspectratio]{#2}%
  }{}%
}

% Explicit placeholder for visuals/examples that will be added later.
\newcommand{\lectureplaceholder}[2][2.35in]{%
  \begin{center}
    \fbox{%
      \begin{minipage}[c][#1][c]{0.90\linewidth}
        \centering
        \small\textit{#2}
      \end{minipage}%
    }
  \end{center}%
}

\setbeamersize{text margin left=0.42in,text margin right=0.42in}

\coursenumber{CS 4880}
\coursetitle{Artificial Intelligence}
\lecturenumber{8}
\lecturetitle{Adversarial Games: Minimax and Expectimax}
\lecturedate{September 15, 2026}
\courseunit{Search and Problem Solving}

\begin{document}

% -----------------------------------------------------------------------------
% Source slide 1 -> New slide 1
% -----------------------------------------------------------------------------
\makedecktitle

% -----------------------------------------------------------------------------
% Source slide 2 -> New slide 2
% -----------------------------------------------------------------------------
\begin{frame}[shrink=18]{Today: Reasoning Against Another Agent}
  \begin{columns}[T,onlytextwidth]
    \begin{column}{0.56\textwidth}
      Early AI treated game playing as a clean test of machine reasoning.
      \begin{itemize}
        \item The rules are explicit and the state is observable.
        \item An agent must reason about possible futures before acting.
        \item The environment contains another decision maker with conflicting goals.
        \item Good decisions depend on predicting how that opponent will respond.
      \end{itemize}
      \begin{block}{Central question}
        How should an agent choose an action when the outcome depends on both its own move and the response of another agent?
      \end{block}
    \end{column}
    \begin{column}{0.38\textwidth}
      \lectureimage[2.55in]{slide2.png}
    \end{column}
  \end{columns}
  \begin{keyidea}
    Adversarial search extends ordinary search by explicitly modeling the choices of other agents.
  \end{keyidea}
\end{frame}

% -----------------------------------------------------------------------------
% Source slide 3 -> New slide 3
% -----------------------------------------------------------------------------
\begin{frame}[shrink=18]{Chess as a Test of Intelligence}
  \begin{columns}[T,onlytextwidth]
    \begin{column}{0.55\textwidth}
      Chess has long been associated with foresight, planning, and strategic reasoning.
      \begin{itemize}
        \item A position admits many legal continuations.
        \item A strong move can only be judged relative to the opponent's best reply.
        \item Success requires looking ahead through alternating choices.
      \end{itemize}
      \begin{block}{Historical motivation}
        Early chess-playing machines helped turn the question ``Can machines think?'' into a concrete computational problem.
      \end{block}
    \end{column}
    \begin{column}{0.39\textwidth}
      \lectureimage[2.65in]{slide3.png}
    \end{column}
  \end{columns}
\end{frame}

% -----------------------------------------------------------------------------
% Source slide 4 -> New slide 4
% -----------------------------------------------------------------------------
\begin{frame}[shrink=18]{Early Computer Chess}
  \begin{columns}[T,onlytextwidth]
    \begin{column}{0.54\textwidth}
      Early work on computer chess established many of the ideas still used in game-playing AI.
      \begin{itemize}
        \item In the 1950s and 1960s, researchers including Alan Turing and Mac Hack explored machine chess.
        \item By the 1970s and 1980s, programs could defeat many amateur players.
        \item The central computational challenge was not learning the rules, but searching enough plausible futures.
      \end{itemize}
      \begin{keyidea}
        Game-playing strength depends heavily on how efficiently the agent searches and evaluates future states.
      \end{keyidea}
    \end{column}
    \begin{column}{0.40\textwidth}
      \lectureimage[2.65in]{slide4.png}
    \end{column}
  \end{columns}
\end{frame}

% -----------------------------------------------------------------------------
% Source slide 5 -> New slide 5
% -----------------------------------------------------------------------------
\begin{frame}[shrink=18]{Deep Blue: First Match}
  \begin{columns}[T,onlytextwidth]
    \begin{column}{0.54\textwidth}
      In 1996, IBM's Deep Blue played a match against world champion Garry Kasparov.
      \begin{itemize}
        \item Deep Blue won the first game of the match.
        \item Kasparov recovered and won the full match $4$--$2$.
        \item The event demonstrated that large-scale search could compete with elite human play in individual games.
      \end{itemize}
      \begin{block}{What mattered computationally?}
        Deep Blue could examine enormous numbers of candidate continuations and evaluate their strategic consequences.
      \end{block}
    \end{column}
    \begin{column}{0.40\textwidth}
      \lectureimage[2.65in]{slide5.png}
    \end{column}
  \end{columns}
\end{frame}

% -----------------------------------------------------------------------------
% Source slide 6 -> New slide 6
% -----------------------------------------------------------------------------
\begin{frame}[shrink=18]{Deep Blue: Rematch}
  \begin{columns}[T,onlytextwidth]
    \begin{column}{0.54\textwidth}
      In the 1997 rematch, Deep Blue defeated Kasparov in a full match under tournament conditions.
      \begin{itemize}
        \item The system relied on specialized hardware and aggressive search.
        \item It evaluated extremely large numbers of chess positions per second.
        \item Its strength came from combining search depth with position evaluation.
      \end{itemize}
      \begin{block}{Transition}
        To understand systems like Deep Blue, we first need a formal model of games and a rule for backing values through a game tree.
      \end{block}
    \end{column}
    \begin{column}{0.40\textwidth}
      \lectureimage[2.65in]{slide6.png}
    \end{column}
  \end{columns}
\end{frame}

\section{Games as Search Problems}

% -----------------------------------------------------------------------------
% Source slide 7 -> New slide 7
% -----------------------------------------------------------------------------
\begin{frame}[shrink=18]{The Classical Adversarial-Game Model}
  We begin with a deliberately restricted class of games.
  \begin{columns}[T,onlytextwidth]
    \begin{column}{0.52\textwidth}
      \begin{itemize}
        \item Two players take turns.
        \item The game is zero-sum: one player's gain is the other's loss.
        \item The state is fully observable.
        \item Actions have deterministic outcomes.
        \item Players act sequentially rather than simultaneously.
      \end{itemize}
      \begin{block}{Examples}
        Tic-tac-toe, checkers, chess, and Go fit much of this model at the rules level.
      \end{block}
    \end{column}
    \begin{column}{0.42\textwidth}
      \lectureimage[2.45in]{slide7.png}
    \end{column}
  \end{columns}
  \begin{keyidea}
    Minimax is most naturally introduced for deterministic, turn-taking, zero-sum, perfect-information games.
  \end{keyidea}
\end{frame}

% -----------------------------------------------------------------------------
% Source slide 8 -> New slide 8
% -----------------------------------------------------------------------------
\begin{frame}[shrink=18]{Formalizing a Game}
  \begin{cpdefinition}
    A deterministic turn-taking game can be described by an initial state, a player-to-move function, legal actions, a transition model, a terminal test, and terminal utilities.
  \end{cpdefinition}
  \begin{columns}[T,onlytextwidth]
    \begin{column}{0.58\textwidth}
      \begin{itemize}
        \item $S_0$: initial game state.
        \item $\mathrm{TO\_MOVE}(s)$: player whose turn it is in state $s$.
        \item $\mathrm{ACTIONS}(s)$: legal actions available in $s$.
        \item $\mathrm{RESULT}(s,a)$: state reached after action $a$.
        \item $\mathrm{IS\_TERMINAL}(s)$: whether the game has ended.
        \item $\mathrm{UTILITY}(s,p)$: payoff to player $p$ in terminal state $s$.
      \end{itemize}
    \end{column}
    \begin{column}{0.36\textwidth}
      \lectureimage[2.35in]{slide8.png}
    \end{column}
  \end{columns}
\end{frame}

% -----------------------------------------------------------------------------
% Source slide 9 -> New slide 9
% -----------------------------------------------------------------------------
\begin{frame}[shrink=18]{Game Trees and Ply}
  \begin{columns}[T,onlytextwidth]
    \begin{column}{0.55\textwidth}
      A \textbf{ply} is one move by one player.
      \begin{itemize}
        \item Nodes represent game states.
        \item Edges represent legal actions.
        \item Levels alternate between the players.
        \item Leaves are terminal states or states where search is intentionally stopped.
        \item Utilities at terminal states describe the final outcome.
      \end{itemize}
      \begin{block}{Perspective}
        In a two-player zero-sum game, we usually measure every terminal value from MAX's perspective.
      \end{block}
    \end{column}
    \begin{column}{0.39\textwidth}
      \lectureimage[2.65in]{slide9.png}
    \end{column}
  \end{columns}
\end{frame}

% -----------------------------------------------------------------------------
% Source slide 10 -> New slide 10
% -----------------------------------------------------------------------------
\begin{frame}[shrink=18]{Utility: What Counts as a Good Outcome?}
  \begin{columns}[T,onlytextwidth]
    \begin{column}{0.53\textwidth}
      For a simple zero-sum game, terminal outcomes can be encoded with a scalar utility.
      \begin{itemize}
        \item MAX win: $+1$.
        \item Draw: $0$.
        \item MAX loss: $-1$.
      \end{itemize}
      \begin{block}{Connect Four example}
        A terminal utility function can inspect the winner and return $+1$, $0$, or $-1$ from MAX's perspective.
      \end{block}
      \begin{keyidea}
        Search does not choose moves directly from labels such as ``win'' or ``loss''; it chooses by comparing utility values.
      \end{keyidea}
    \end{column}
    \begin{column}{0.41\textwidth}
      \lectureimage[2.60in]{slide10.png}
    \end{column}
  \end{columns}
\end{frame}

% -----------------------------------------------------------------------------
% Source slide 11 -> New slide 11
% -----------------------------------------------------------------------------
\begin{frame}[shrink=18]{Why Ordinary Search Is Not Enough}
  \begin{columns}[T,onlytextwidth]
    \begin{column}{0.55\textwidth}
      Classical search assumes the environment does not deliberately interfere with the plan.
      \begin{itemize}
        \item In a game, the next state is partly determined by another agent.
        \item A plan that succeeds against one response may fail against another.
        \item The opponent is not merely an obstacle; it is a decision maker with its own objective.
      \end{itemize}
      \begin{block}{Questions we must model}
        Will the opponent choose its strongest move? Will it behave randomly? Is some uncertainty outside either player's control?
      \end{block}
    \end{column}
    \begin{column}{0.39\textwidth}
      \lectureimage[2.55in]{slide11.png}
    \end{column}
  \end{columns}
\end{frame}

% -----------------------------------------------------------------------------
% Source slide 12 -> New slide 12
% -----------------------------------------------------------------------------
\begin{frame}[shrink=18]{Decision Nodes: MAX, MIN, and Chance}
  \begin{columns}[T,onlytextwidth]
    \begin{column}{0.52\textwidth}
      Different node types represent different causes of the next state.
      \begin{itemize}
        \item \textbf{MAX node}: our agent chooses an action.
        \item \textbf{MIN node}: an adversarial opponent chooses an action.
        \item \textbf{Chance node}: an outcome is sampled according to a probability distribution.
      \end{itemize}
      \begin{block}{Three backup rules}
        Maximum, minimum, and expectation are three different ways to summarize the values of child states.
      \end{block}
    \end{column}
    \begin{column}{0.42\textwidth}
      \lectureimage[2.55in]{slide12.png}
    \end{column}
  \end{columns}
  \begin{keyidea}
    The structure of the environment determines how values should be propagated upward through the tree.
  \end{keyidea}
\end{frame}

% -----------------------------------------------------------------------------
% Source slide 13 -> New slide 13
% -----------------------------------------------------------------------------
\begin{frame}[shrink=18]{Tic-Tac-Toe as a Game Tree}
  \begin{columns}[T,onlytextwidth]
    \begin{column}{0.54\textwidth}
      Suppose MAX plays X and moves first.
      \begin{itemize}
        \item From the empty board, MAX has nine legal moves.
        \item After each X move, MIN chooses an O move.
        \item Play alternates until a win, loss, or draw occurs.
        \item Terminal boards receive utilities such as $+1$, $0$, and $-1$.
      \end{itemize}
      \begin{block}{Tree interpretation}
        Each complete root-to-leaf path is one possible play sequence; internal values summarize what rational play implies from that state.
      \end{block}
    \end{column}
    \begin{column}{0.40\textwidth}
      \lectureimage[2.70in]{slide13.png}
    \end{column}
  \end{columns}
\end{frame}

\section{Minimax}

% -----------------------------------------------------------------------------
% Source slide 14 -> New slide 14
% -----------------------------------------------------------------------------
\begin{frame}[shrink=18]{Minimax}
  \begin{cpdefinition}
    \textbf{Minimax} assigns a value to each game state assuming MAX chooses actions that maximize utility and MIN chooses actions that minimize MAX's utility.
  \end{cpdefinition}
  \begin{columns}[T,onlytextwidth]
    \begin{column}{0.52\textwidth}
      \begin{itemize}
        \item MAX nodes choose the child with the largest value.
        \item MIN nodes choose the child with the smallest value.
        \item Terminal states supply the base-case utilities.
        \item Values are backed up recursively from leaves toward the root.
      \end{itemize}
    \end{column}
    \begin{column}{0.42\textwidth}
      \lectureimage[2.50in]{slide14.png}
    \end{column}
  \end{columns}
  \begin{keyidea}
    Minimax asks: which move gives MAX the best outcome that can still be guaranteed against an optimal opponent?
  \end{keyidea}
\end{frame}

% -----------------------------------------------------------------------------
% Source slide 15 -> New slide 15
% -----------------------------------------------------------------------------
\begin{frame}[shrink=18]{Backing Up Minimax Values}
  \begin{columns}[T,onlytextwidth]
    \begin{column}{0.56\textwidth}
      For a state $s$:
      \[
        V(s)=
        \begin{cases}
          \mathrm{UTILITY}(s), & \text{if $s$ is terminal},\\
          \max\limits_{a\in \mathrm{ACTIONS}(s)} V(\mathrm{RESULT}(s,a)), & \text{if MAX moves},\\
          \min\limits_{a\in \mathrm{ACTIONS}(s)} V(\mathrm{RESULT}(s,a)), & \text{if MIN moves}.
        \end{cases}
      \]
      \begin{itemize}
        \item MIN assumes MAX will exploit every weakness.
        \item MAX assumes MIN will choose the most damaging reply.
      \end{itemize}
    \end{column}
    \begin{column}{0.38\textwidth}
      \lectureimage[2.45in]{slide15.png}
    \end{column}
  \end{columns}
\end{frame}

% -----------------------------------------------------------------------------
% Source slide 16 -> New slide 16
% -----------------------------------------------------------------------------
\begin{frame}[shrink=18]{Worked Example: A Two-Ply Minimax Tree}
  \begin{columns}[T,onlytextwidth]
    \begin{column}{0.52\textwidth}
      Consider a root MAX node with three legal actions.
      \begin{itemize}
        \item Each action leads to a MIN node.
        \item Each MIN node has several terminal children.
        \item MIN first selects the smallest child value in each subtree.
        \item MAX then chooses the largest of those backed-up values.
      \end{itemize}
      \begin{block}{In-class calculation}
        Fill in the MIN values first, then determine the action selected at the root.
      \end{block}
    \end{column}
    \begin{column}{0.42\textwidth}
      \lectureimage[2.70in]{slide16.png}
    \end{column}
  \end{columns}
  \begin{keyidea}
    The root value is the utility MAX can force when both players choose optimally.
  \end{keyidea}
\end{frame}

% -----------------------------------------------------------------------------
% Source slide 17 -> New slide 17
% -----------------------------------------------------------------------------
\begin{frame}[shrink=18]{The Minimax Search Algorithm}
  \begin{columns}[T,onlytextwidth]
    \begin{column}{0.58\textwidth}
      \begin{block}{Recursive structure}
      \ttfamily\scriptsize
      MINIMAX(state):\\[0.03in]
      \quad if TERMINAL(state):\\
      \qquad return UTILITY(state)\\[0.03in]
      \quad if TO\_MOVE(state) = MAX:\\
      \qquad value $\leftarrow -\infty$\\
      \qquad for action in ACTIONS(state):\\
      \qquad\quad value $\leftarrow \max$(value, MINIMAX(RESULT(state,action)))\\
      \qquad return value\\[0.03in]
      \quad else:\\
      \qquad value $\leftarrow +\infty$\\
      \qquad for action in ACTIONS(state):\\
      \qquad\quad value $\leftarrow \min$(value, MINIMAX(RESULT(state,action)))\\
      \qquad return value
      \end{block}
    \end{column}
    \begin{column}{0.36\textwidth}
      \begin{itemize}
        \item Depth-first recursion explores the game tree.
        \item Terminal utilities provide the base case.
        \item Internal nodes aggregate child values by max or min.
      \end{itemize}
      \lectureplaceholder[1.55in]{Code/example placeholder: Python implementation using a GameState interface.}
    \end{column}
  \end{columns}
\end{frame}

% -----------------------------------------------------------------------------
% Source slide 18 -> New slide 18
% -----------------------------------------------------------------------------
\begin{frame}[shrink=18]{Minimax Properties}
  \begin{columns}[T,onlytextwidth]
    \begin{column}{0.48\textwidth}
      \begin{block}{Correctness}
        \begin{itemize}
          \item Complete for finite game trees.
          \item Optimal against an optimal opponent.
        \end{itemize}
      \end{block}
      \begin{block}{Time}
        With branching factor $b$ and depth $m$:
        \[
          O(b^m).
        \]
      \end{block}
    \end{column}
    \begin{column}{0.48\textwidth}
      \begin{block}{Space}
        Depth-first recursive implementations require memory proportional to the search depth plus generated successors.
      \end{block}
      \begin{block}{Practical consequence}
        The exponential number of nodes, not the minimax backup itself, is the central computational bottleneck.
      \end{block}
      \lectureplaceholder[1.20in]{Visual placeholder: small game tree annotated with $b$ and $m$.}
    \end{column}
  \end{columns}
  \begin{keyidea}
    Exact minimax is conceptually simple but quickly becomes infeasible as the game tree grows.
  \end{keyidea}
\end{frame}

% -----------------------------------------------------------------------------
% Source slide 19 -> New slide 19
% -----------------------------------------------------------------------------
\begin{frame}[shrink=18]{Historical Roots: Minimax and Game Theory}
  \begin{columns}[T,onlytextwidth]
    \begin{column}{0.48\textwidth}
      \begin{block}{John von Neumann}
        The minimax theorem established a foundation for rational play in two-player zero-sum games.
      \end{block}
      \begin{itemize}
        \item Formalized adversarial utility.
        \item Motivated worst-case decision making.
        \item Connected strategic choice to mathematical game theory.
      \end{itemize}
    \end{column}
    \begin{column}{0.48\textwidth}
      \begin{block}{Beyond zero-sum games}
        Later game theory, including Nash equilibrium, studies settings in which utilities are not exact opposites and stable strategic outcomes can involve multiple agents.
      \end{block}
      \lectureimage[1.75in]{slide19.png}
    \end{column}
  \end{columns}
\end{frame}

\section{Beyond Two-Player Zero-Sum Games}

% -----------------------------------------------------------------------------
% Source slide 20 -> New slide 20
% -----------------------------------------------------------------------------
\begin{frame}[shrink=18]{Non-Zero-Sum Games}
  \begin{columns}[T,onlytextwidth]
    \begin{column}{0.55\textwidth}
      In non-zero-sum games, players' utilities are not exact negatives of one another.
      \begin{itemize}
        \item Terminal states may return a utility vector such as $(u_A,u_B)$.
        \item Both players can benefit from some outcomes.
        \item One player's improvement need not imply an equal loss for another.
      \end{itemize}
      \begin{block}{Consequence}
        A single scalar ``MAX utility'' is no longer sufficient to describe every player's preference.
      \end{block}
    \end{column}
    \begin{column}{0.39\textwidth}
      \lectureimage[2.60in]{slide20.png}
    \end{column}
  \end{columns}
\end{frame}

% -----------------------------------------------------------------------------
% Source slide 21 -> New slide 21
% -----------------------------------------------------------------------------
\begin{frame}[shrink=18]{Cooperation in Non-Zero-Sum Games}
  \begin{columns}[T,onlytextwidth]
    \begin{column}{0.54\textwidth}
      When utilities are not exact opposites, cooperation can be rational.
      \begin{itemize}
        \item A joint action may improve the payoff of both players.
        \item Competition and cooperation can coexist in the same game.
        \item Strategy depends on each player's own utility, not on a universal win/loss score.
      \end{itemize}
      \begin{block}{Illustration}
        A cooperative outcome such as $(1000,1000)$ may dominate an aggressive outcome such as $(1000,0)$ for one participant without being zero-sum.
      \end{block}
    \end{column}
    \begin{column}{0.40\textwidth}
      \lectureimage[2.55in]{slide21.png}
    \end{column}
  \end{columns}
\end{frame}

% -----------------------------------------------------------------------------
% Source slide 22 -> New slide 22
% -----------------------------------------------------------------------------
\begin{frame}[shrink=18]{Extending Search to Multiple Players}
  \begin{columns}[T,onlytextwidth]
    \begin{column}{0.53\textwidth}
      With three or more players, utilities naturally become vectors.
      \begin{itemize}
        \item Example terminal utility: $(u_A,u_B,u_C)$.
        \item Each player evaluates outcomes according to its own component.
        \item The player to move chooses the successor that maximizes its own utility component.
      \end{itemize}
      \begin{block}{Generalization}
        This is often called \textbf{max$^n$}: each node backs up the utility vector of the child preferred by the current player.
      \end{block}
    \end{column}
    \begin{column}{0.41\textwidth}
      \lectureimage[2.65in]{slide22.png}
    \end{column}
  \end{columns}
\end{frame}

% -----------------------------------------------------------------------------
% Source slide 23 -> New slide 23
% -----------------------------------------------------------------------------
\begin{frame}[shrink=18]{Choosing Moves in Multiplayer Search}
  \begin{columns}[T,onlytextwidth]
    \begin{column}{0.54\textwidth}
      Suppose player C chooses between terminal vectors $(1,2,6)$ and $(4,2,3)$.
      \begin{itemize}
        \item C compares the third component: $6$ versus $3$.
        \item C therefore chooses $(1,2,6)$.
        \item The entire vector is then backed up to C's parent.
      \end{itemize}
      \begin{block}{Important distinction}
        Each player maximizes a different component of the same backed-up vector.
      \end{block}
    \end{column}
    \begin{column}{0.40\textwidth}
      \lectureimage[2.55in]{slide23.png}
    \end{column}
  \end{columns}
\end{frame}

% -----------------------------------------------------------------------------
% Source slide 24 -> New slide 24
% -----------------------------------------------------------------------------
\begin{frame}[shrink=18]{Alliances and Emergent Cooperation}
  \begin{columns}[T,onlytextwidth]
    \begin{column}{0.54\textwidth}
      Multiplayer games can produce strategic behavior that does not appear in strict two-player minimax.
      \begin{itemize}
        \item Two weak players may both prefer to attack a stronger third player.
        \item Their actions can align even without explicit coordination.
        \item Apparent cooperation may emerge from individual utility maximization.
      \end{itemize}
      \begin{block}{Game-tree explanation}
        If attacking the stronger player improves each participant's own payoff, both can independently choose the same strategic direction.
      \end{block}
    \end{column}
    \begin{column}{0.40\textwidth}
      \lectureimage[2.55in]{slide24.png}
    \end{column}
  \end{columns}
\end{frame}

\section{Expectimax and Chance}

% -----------------------------------------------------------------------------
% Source slide 25 -> New slide 25
% -----------------------------------------------------------------------------
\begin{frame}[shrink=18]{When the Next Move Is Not Adversarial}
  \begin{columns}[T,onlytextwidth]
    \begin{column}{0.47\textwidth}
      \begin{block}{Minimax assumption}
        \begin{itemize}
          \item The other agent chooses deliberately.
          \item It selects the action that minimizes our utility.
          \item We plan against its strongest response.
        \end{itemize}
      \end{block}
    \end{column}
    \begin{column}{0.47\textwidth}
      \begin{block}{Stochastic assumption}
        \begin{itemize}
          \item The next outcome may be random or noisy.
          \item Different outcomes occur with known or modeled probabilities.
          \item We should optimize expected utility rather than worst-case utility.
        \end{itemize}
      \end{block}
    \end{column}
  \end{columns}
  \begin{center}
    \lectureimage[1.45in]{slide25.png}
  \end{center}
  \begin{keyidea}
    Minimax uses a minimum backup for an adversary; expectimax uses an expectation backup for chance.
  \end{keyidea}
\end{frame}

% -----------------------------------------------------------------------------
% Source slide 26 -> New slide 26
% -----------------------------------------------------------------------------
\begin{frame}[shrink=18]{Expectimax}
  \begin{cpdefinition}
    \textbf{Expectimax} evaluates decision trees containing MAX nodes and chance nodes by maximizing at decision nodes and taking probability-weighted averages at chance nodes.
  \end{cpdefinition}
  \begin{columns}[T,onlytextwidth]
    \begin{column}{0.55\textwidth}
      \begin{itemize}
        \item At a MAX node: choose the child with maximum value.
        \item At a chance node: compute the expected value of the children.
      \end{itemize}
      If chance node $s$ has outcomes $s_1,\ldots,s_k$ with probabilities $p_1,\ldots,p_k$:
      \[
        V(s)=\sum_{i=1}^{k} p_i V(s_i).
      \]
      \begin{block}{Expectiminimax}
        Trees that contain MAX, MIN, \emph{and} chance nodes combine all three backup rules.
      \end{block}
    \end{column}
    \begin{column}{0.39\textwidth}
      \lectureimage[2.20in]{slide26.png}
    \end{column}
  \end{columns}
\end{frame}

% -----------------------------------------------------------------------------
% Source slide 27 -> New slide 27
% -----------------------------------------------------------------------------
\begin{frame}[shrink=18]{Worked Example: Expected Utility Backup}
  \begin{columns}[T,onlytextwidth]
    \begin{column}{0.54\textwidth}
      Suppose a chance node has two equally likely outcomes with backed-up values $3$ and $2$.
      \[
        V(\text{chance})=0.5(3)+0.5(2)=2.5.
      \]
      \begin{itemize}
        \item The chance node does not select either child.
        \item It represents the average utility induced by the probability model.
        \item A MAX parent compares this expected value against the expected values of alternative actions.
      \end{itemize}
      \begin{block}{In-class extension}
        Change the probabilities to $0.8$ and $0.2$. Does MAX's preferred root action change?
      \end{block}
    \end{column}
    \begin{column}{0.40\textwidth}
      \lectureimage[2.65in]{slide27.png}
    \end{column}
  \end{columns}
\end{frame}

% -----------------------------------------------------------------------------
% Source slide 28 -> New slide 28
% -----------------------------------------------------------------------------
\begin{frame}[shrink=18]{Adversarial Search in Popular Culture: WarGames}
  \begin{columns}[T,onlytextwidth]
    \begin{column}{0.54\textwidth}
      \emph{WarGames} (1983) dramatizes a computer reasoning about competitive strategic outcomes.
      \begin{itemize}
        \item The film treats games as models of strategic interaction.
        \item The famous conclusion is that some adversarial settings have no desirable winning continuation.
        \item This provides a useful discussion point: utility design determines what the search regards as a successful outcome.
      \end{itemize}
      \begin{block}{Discussion prompt}
        If every terminal outcome has catastrophic utility, what should a rational decision procedure prefer?
      \end{block}
    \end{column}
    \begin{column}{0.40\textwidth}
      \lectureimage[2.55in]{slide28.png}
    \end{column}
  \end{columns}
\end{frame}

% -----------------------------------------------------------------------------
% Source slide 29 -> New slide 29
% -----------------------------------------------------------------------------
\begin{frame}[shrink=18]{What If Every Branch Is Bad?}
  \begin{columns}[T,onlytextwidth]
    \begin{column}{0.54\textwidth}
      Search still behaves rationally when all outcomes are undesirable.
      \begin{itemize}
        \item Minimax selects the action with the best worst-case utility.
        \item Expectimax selects the action with the highest expected utility.
        \item Neither algorithm requires the existence of a positive or ``winning'' terminal state.
      \end{itemize}
      \begin{block}{Important consequence}
        A rational search agent may choose the least harmful outcome rather than a true victory.
      \end{block}
    \end{column}
    \begin{column}{0.40\textwidth}
      \lectureimage[2.55in]{slide29.png}
    \end{column}
  \end{columns}
  \begin{keyidea}
    Utility defines preference, not success in an absolute sense.
  \end{keyidea}
\end{frame}

% -----------------------------------------------------------------------------
% Source slide 30 -> New slide 30
% -----------------------------------------------------------------------------
\begin{frame}[shrink=18]{Adversarial Search in Video Games}
  \begin{columns}[T,onlytextwidth]
    \begin{column}{0.54\textwidth}
      Game-tree reasoning appears naturally in turn-based and tactical games.
      \begin{itemize}
        \item Strategy games: Civilization, XCOM.
        \item Board-game adaptations: chess, checkers, Go.
        \item Tactical combat: Fire Emblem, Advance Wars.
      \end{itemize}
      A game-playing agent can:
      \begin{itemize}
        \item generate candidate actions and counteractions;
        \item evaluate future states;
        \item back values toward the current state;
        \item select the action with the strongest backed-up value.
      \end{itemize}
    \end{column}
    \begin{column}{0.40\textwidth}
      \lectureplaceholder[2.65in]{Visual placeholder: turn-based strategy screenshot plus a small game tree.}
    \end{column}
  \end{columns}
\end{frame}

% -----------------------------------------------------------------------------
% Source slide 31 -> New slide 31
% -----------------------------------------------------------------------------
\begin{frame}[shrink=18]{Chess Moves Become a Search Tree}
  \begin{columns}[T,onlytextwidth]
    \begin{column}{0.48\textwidth}
      From one chess position:
      \begin{itemize}
        \item enumerate candidate legal moves;
        \item generate resulting positions;
        \item alternate the side to move;
        \item continue recursively.
      \end{itemize}
      \begin{block}{Branching}
        Each additional ply multiplies the number of positions that may need to be considered.
      \end{block}
    \end{column}
    \begin{column}{0.46\textwidth}
      \lectureplaceholder[2.70in]{Example placeholder: chessboard with two candidate moves beside the corresponding first levels of a search tree.}
    \end{column}
  \end{columns}
\end{frame}

% -----------------------------------------------------------------------------
% Source slide 32 -> New slide 32
% -----------------------------------------------------------------------------
\begin{frame}[shrink=18]{Chess Trees Are Too Large for Terminal Search}
  \begin{columns}[T,onlytextwidth]
    \begin{column}{0.54\textwidth}
      Full minimax assumes that search eventually reaches terminal outcomes.
      \begin{itemize}
        \item In chess, the game tree is far too large to search to completion from typical positions.
        \item Practical systems therefore stop at a limited depth.
        \item Nonterminal frontier states must be scored by an \textbf{evaluation function} rather than a true terminal utility.
      \end{itemize}
      \begin{block}{Evaluation idea}
        Estimate position quality using features such as material, king safety, mobility, or positional structure.
      \end{block}
    \end{column}
    \begin{column}{0.40\textwidth}
      \lectureplaceholder[2.60in]{Example placeholder: chess position with a heuristic evaluation score and a depth-limited tree.}
    \end{column}
  \end{columns}
  \begin{keyidea}
    Practical adversarial search replaces unreachable terminal utilities with heuristic estimates at a cutoff depth.
  \end{keyidea}
\end{frame}

% -----------------------------------------------------------------------------
% Source slide 33 -> New slide 33
% -----------------------------------------------------------------------------
\begin{frame}[shrink=18]{Deep Blue and the Cost of Search}
  \begin{columns}[T,onlytextwidth]
    \begin{column}{0.54\textwidth}
      Deep Blue's strength depended on searching huge numbers of positions efficiently.
      \begin{itemize}
        \item Game trees grow exponentially with depth.
        \item Much of the tree may be irrelevant to the final decision.
        \item Searching deeper is valuable only if the computation can be focused effectively.
      \end{itemize}
      \begin{block}{Next question}
        Can we avoid exploring branches that cannot possibly change the minimax decision?
      \end{block}
    \end{column}
    \begin{column}{0.40\textwidth}
      \lectureplaceholder[2.60in]{Visual placeholder: Deep Blue image plus an exponentially expanding tree with several branches faded out.}
    \end{column}
  \end{columns}
\end{frame}

% -----------------------------------------------------------------------------
% Source slide 34 -> New slide 34
% -----------------------------------------------------------------------------
\begin{frame}[shrink=18]{Takeaways}
  \begin{itemize}
    \item Adversarial search models environments in which another agent actively chooses actions that affect our outcome.
    \item Minimax assumes a rational opponent and backs values up using alternating maximum and minimum operations.
    \item Exact minimax is optimal for finite deterministic zero-sum games but requires exponential time in the search depth.
    \item Multiplayer and non-zero-sum games require richer utility representations, often utility vectors.
    \item Expectimax replaces an adversarial minimum with a probability-weighted expectation when uncertainty is stochastic rather than strategic.
    \item Practical game-playing systems use depth limits and evaluation functions because real game trees are too large to solve exhaustively.
    \item Next lecture: \textbf{Alpha--Beta Pruning}.
  \end{itemize}
  \begin{keyidea}
    Minimax plans against the strongest opponent response; expectimax plans against a distribution over uncertain outcomes.
  \end{keyidea}
\end{frame}

\end{document}
